\documentclass{article}

\usepackage[preprint]{tmlr}

\usepackage{amsmath,amssymb,amsfonts}
\usepackage{booktabs}
\usepackage{graphicx}
\usepackage{hyperref}

\renewcommand{\arraystretch}{1.3}
\setlength{\abovetopsep}{6pt}

\title{Bracket Norm Is Not Mechanism:\\ Validity Conditions for Noncommutative Perturbation Scores}

\author{\name Elliot Tower \\
      \email elliot@elliottower.ai}

\begin{document}
\maketitle

\begin{abstract}
Direction instability---the mean pairwise angular deviation of perturbation
signatures across contexts---predicts drug mechanism transport across cell
lines. We ask what this prediction \emph{licenses}. Through pre-registered
experiments across three perturbation domains---8{,}949 LINCS L1000
compounds (978 transcriptomic features), 1{,}676 Perturb-seq gene
knockdowns \citep{replogle2022mapping} (2{,}073 genes), and 25{,}254
JUMP-CP compounds (3{,}180 Cell Painting morphological features)---we show
that validity conditions do not reclassify perturbations: they determine
what claims a practitioner can make about the ranking. In all three domains,
correction preserves the ranking (LINCS: $\rho > 0.99$; Perturb-seq: $\rho
= 0.91$; JUMP-CP: $\rho = 0.99$). Three drug-domain validity conditions
are tested empirically: phenotype-projected bracket correlates with
on-target genetic perturbation connectivity ($\rho = 0.38$, $p < 10^{-27}$)
while raw bracket does not ($\rho = -0.04$); transport-stable bracket
outpredicts raw bracket for held-out cell-line effects in all 5
cross-validation folds; and localization to target-gene knockdown signatures
improves prediction for intracellular enzyme inhibitors but not
receptor-mediated drugs, producing an informative null. The validity
ladder functions as a proof system for mechanistic claims: each rung does not
change the score, it upgrades the inferential warrant attached to that score.
\end{abstract}

\section{Introduction}

The Connectivity Map paradigm \citep{lamb2006connectivity,subramanian2017next}
established that transcriptomic-signature similarity predicts drug mechanism
of action. Direction instability extends this from pairwise similarity to
angular consistency: rather than asking whether two drugs produce similar
signatures, it asks whether a single drug's signature rotates across cell
lines. \citet{tower2026bracket} applied this metric to 8{,}949 drugs in
LINCS L1000, showing that mechanism-conserving compounds (pan-HDAC
inhibitors, topoisomerase inhibitors, CDK inhibitors) achieve low direction
instability while receptor-mediated drugs cluster near the population mean
of 0.92. The HDAC selectivity gradient---pan-isoform
inhibitors at $D = 0.56$--$0.63$, selective inhibitors at $D = 0.96$--$0.99$---provides
a natural controlled comparison within a single drug class.

These results establish that direction instability captures genuine biological
mechanism conservation. They leave open a different question: under what
conditions does a direction-instability score license a mechanistic \emph{claim}?
A score of $D = 0.55$ means signatures point in similar directions across cell
lines. It does not, by itself, mean the drug acts through a specific pathway,
that the effect is therapeutic rather than cytotoxic, or that the observation
will replicate in a new dataset.

Confound corrections for perturbation scores are well-established:
CERES corrects copy-number artifacts in CRISPR screens
\citep{meyers2017ceres}; connectivity-metric evaluations benchmark
signature-similarity scores \citep{li2020comprehensive}. These
corrections improve the \emph{accuracy} of the score. We ask a different
question: when does a score---even an accurate one---license a mechanistic
\emph{claim}? A correction that preserves the global ranking ($\rho > 0.99$)
adds no discriminative power, yet it upgrades the inferential warrant by
ruling out a specific confounding explanation.

We formalize this through a validity ladder: a sequence of conditions, each
building on the previous, that progressively strengthen what a
direction-instability measurement licenses. The ladder has seven rungs, from
declaring the process view (what counts as a context?) through phenotype
alignment and cross-context transport to holdout prediction. Each rung
answers a specific threat to interpretation. The ladder is a proof system,
analogous to the Bradford Hill criteria for causal inference and to recent
construct-validity frameworks for ML benchmarks
\citep{freiesleben2025benchmarking}: satisfying more rungs does not change
the number, it changes what you are allowed to conclude from it.

This paper tests whether validity conditions matter empirically across three
perturbation domains. In LINCS L1000, pre-registered experiments probe
specific failure modes of drug direction instability: cytotoxicity
masquerading as mechanism conservation and broad therapeutics that are
smooth rather than mechanism-absent. In Perturb-seq genetic knockdowns from
Replogle et al.\ \citep{replogle2022mapping}, we test whether essential-gene
correction---the genetic analogue of toxicity correction---changes the
Grassmannian geodesic distance ranking between K562 and RPE1 perturbation
response subspaces. In JUMP-CP Cell Painting morphological profiles
\citep{chandrasekaran2023jump}, we test whether removing cell-health
features---the morphological analogue of toxicity and essentiality
correction---changes the direction-instability ranking across 25{,}254
compounds measured in $\geq$5 imaging sources. We report honest nulls
alongside confirmations, because the pattern of where validity conditions
bite---and where they do not---is itself the finding.

\section{Background}

\subsection{Direction instability}

For a drug tested in $K$ cell lines, let $\mathbf{s}_1, \ldots,
\mathbf{s}_K \in \mathbb{R}^{978}$ denote its consensus signature vectors.
Define unit-normalized signatures $\hat{\mathbf{s}}_i = \mathbf{s}_i /
\|\mathbf{s}_i\|_2$. Direction instability is:
\begin{equation}
  D = 1 - \frac{2}{K(K-1)} \sum_{i < j}
    \hat{\mathbf{s}}_i^\top \hat{\mathbf{s}}_j
  \label{eq:direction_instability}
\end{equation}
$D = 0$ when all signatures are collinear. $D = 1$ when orthogonal on average.
(The term ``bracket norm'' in the companion paper
\citep{tower2026bracket} refers to this same quantity viewed as a
noncommutative perturbation score; we use ``direction instability''
throughout to avoid terminological collision with unrelated norm
constructions.)
This metric was introduced in \citet{tower2026bracket} to characterize drug mechanism
transport and was shown to anticorrelate with gene-level Jaccard overlap of
top perturbed genes (Spearman $\rho = -0.79$).

\subsection{Grassmannian geodesic distance in Perturb-seq}

Direction instability generalizes from vectors to subspaces. For a gene
knocked down in two cell types, let $U_1, U_2 \in \text{Gr}(k, d)$ denote
the top-$k$ PCA subspaces of the perturbation response in each cell type.
The Grassmannian geodesic distance $\delta(U_1, U_2) =
\|\boldsymbol{\theta}\|_2$, where $\boldsymbol{\theta}$ are the principal
angles between the subspaces, measures how much the perturbation response
\emph{rotates} across cell types. With $K = 2$ contexts, this reduces to a
single pairwise comparison---the subspace analogue of direction instability
for a drug tested in two cell lines.

\subsection{Direction instability in Cell Painting morphological profiles}

Cell Painting \citep{chandrasekaran2023jump} measures 3{,}180 morphological
features per well (cell area, shape, granularity, texture, and intensity
across five fluorescent channels: DNA, RNA, ER, AGP, Mito). JUMP-CP
profiles 116{,}000 compounds across 10 imaging sources. Direction
instability applies directly: each compound's morphological signature in a
given source is a vector in $\mathbb{R}^{3180}$, and $D$ measures angular
consistency across sources. A subset of features---cell area, compactness,
eccentricity, DNA/Mito intensity, granularity, and radial
distribution---reflect general cell health rather than compound-specific
morphological effects. The cell-health correction removes these 743 features
before recomputing direction instability, analogous to the toxicity
correction in LINCS and the essential-gene correction in Perturb-seq.

\subsection{The validity problem}

Consider two drugs with identical direction instability, $D = 0.55$:

\begin{itemize}
  \item \textbf{Drug A} is a pan-HDAC inhibitor. Its conserved signature
    reflects chromatin remodeling genes (\emph{SUV39H1}, \emph{MYC},
    \emph{CDK6}). The consistency comes from targeting a universally
    expressed enzymatic family.
  \item \textbf{Drug B} is a mitochondrial toxin. Its conserved signature
    reflects a generic stress response (\emph{HSPA1A}, \emph{DDIT3},
    \emph{ATF4}). The consistency comes from triggering the same damage
    pathway regardless of cell identity.
\end{itemize}

Both achieve $D = 0.55$. The interpretation differs entirely: Drug A
transports a therapeutic mechanism; Drug B transports cellular violence. The
raw metric cannot distinguish these cases. A validity condition---toxicity
correction, in this instance---can.

\subsection{The validity ladder}

We define seven rungs:

\begin{enumerate}
  \item \textbf{Process view declared.} What constitutes a ``context''?
    (Cell lines, cell types, patients, time points.)
  \item \textbf{Reliable vector field.} Signatures have adequate signal-to-noise
    for direction to be meaningful.
  \item \textbf{Perturbation estimable.} The effect can be isolated from
    background variation.
  \item \textbf{Localized to biology.} Instability concentrates in
    pathway-relevant genes, not globally.
  \item \textbf{Phenotype-aligned.} The direction of conservation matches a
    known target pathway.
  \item \textbf{Transports across contexts.} The score replicates in held-out
    contexts (low Fr\'echet variance).
  \item \textbf{Predicts holdout.} The metric predicts an external outcome
    in new data.
\end{enumerate}

Rungs 1--3 are prerequisites for any directional metric. Rungs 4--7 are the
validity conditions we test empirically: each eliminates a specific
interpretive threat without altering the underlying score.

\section{Related work}

\paragraph{Confound correction in genetic screens.}
Our essential-gene correction (Experiment~7) addresses a problem with
extensive prior art in dependency-screen analysis. CERES
\citep{meyers2017ceres} corrects copy-number-induced false positives in
CRISPR screens by modeling gene-independent effects; integrated
Broad--Sanger analyses \citep{pacini2021integrated} identify 2{,}103 common
essentials and document cross-study divergence of the kind we observe between
K562 and RPE1. These corrections operate on dependency-score magnitudes. Our
contribution is orthogonal: we correct subspace geometry (Grassmannian
geodesic distance), showing that corrections are inferentially meaningful
even when they are rank-preserving at the population level.

\paragraph{Connectivity-based MOA discovery.}
The Connectivity Map framework \citep{subramanian2017next,lamb2006connectivity}
established that transcriptomic-signature similarity predicts mechanism of
action. Subsequent work evaluated connectivity metrics systematically
\citep{li2020comprehensive}. Direction instability extends this paradigm from
pairwise signature similarity to angular consistency across contexts: rather
than asking whether two drugs produce similar signatures, we ask whether a
single drug's signature rotates across cell lines.

\paragraph{Subspace geometry on perturbation data.}
Grassmannian methods for single-cell perturbation analysis are an active
area. LEMUR \citep{ahlmanneltze2025lemur} performs multivariate regression on
Grassmann manifolds to disentangle sources of variance in multi-condition
Perturb-seq data. Our Experiment~7 uses a simpler variant---Grassmannian
geodesic distance between per-gene PCA subspaces across two cell types---but
the geometric language is shared.

\paragraph{Construct validity for computational scores.}
The validity ladder draws on a broader conversation about when benchmark
scores license scientific inferences.
\citet{freiesleben2025benchmarking} develop construct-validity conditions
inspired by psychological measurement theory for evaluating ML models;
\citet{watson2025measuring} propose an eight-recommendation checklist for
LLM benchmarks mapping onto recognized validity dimensions.
Our ladder is the biological-perturbation counterpart: each rung corresponds
to a validity dimension (define the phenomenon, measure only the phenomenon,
analyze confounders, predict external outcomes). In the perturbation-prediction
setting, the SBB (Signal, Bounds, and Baselines) framework
\citep{sbb2026principles} evaluates whether a model's predictions are
trustworthy; our ladder evaluates whether a metric's score licenses a
mechanistic claim.

\section{Methods}

\subsection{Data}

Perturbation signatures were drawn from LINCS L1000 Phase I (GEO accession
GSE92742; $\sim$1.3 million profiles across $\sim$50 cell lines and
$\sim$20{,}000 compounds), using Level 5 replicate-collapsed z-scores
across 978 landmark genes \citep{subramanian2017next}. We retained 8{,}949
small-molecule perturbagens tested in $\geq5$ cell lines, the same dataset
characterized in \citet{tower2026bracket}.

Perturb-seq data derive from Replogle et al.\
\citep{replogle2022mapping}: genome-scale CRISPRi knockdowns in K562
(chronic myelogenous leukemia) and RPE1 (retinal pigment epithelium)
cells. We extracted 1{,}676 genes with $\geq$30 perturbed cells in both
cell types, computed per-gene expression responses over 2{,}073 highly
variable genes, and fitted $k = 5$ PCA subspaces per gene per cell type.
Grassmannian geodesic distances were computed between the K562 and RPE1
subspaces for each gene. Both the K562 and RPE1 screens use guide libraries
targeting DepMap-essential genes \citep{replogle2022mapping}; 85\% of RPE1
targets overlap with K562-essential targets (1{,}875 of 2{,}204). The
``non-essential'' comparison group therefore consists of genes that were
selected for essentiality but fall below the stringent Chronos $< -0.5$
threshold in one or both cell types, making this a conservative test of
essential vs.\ non-essential divergence. Essentiality labels derive from
DepMap Chronos scores \citep{dempster2021chronos}: 855 genes with Chronos
$< -0.5$ in both K562 and RPE1 were classified as universal-essential.

JUMP-CP Cell Painting profiles derive from the Cell Painting Gallery
\citep{chandrasekaran2023jump}: variance-stabilized, MAD-normalized
compound profiles across 10 imaging sources. The full dataset contains
116{,}750 compound profiles; after deduplication and filtering to compounds
present in $\geq$5 sources, 25{,}254 compounds remain (78\% attrition,
primarily from compounds tested in $<$5 sources). We aggregated per
compound $\times$ source by mean and computed direction instability over
the full 3{,}180 feature set. Cell-health features (743 features matching Area, Compactness,
Eccentricity, DNA/Mito Intensity, Granularity, Radial Distribution, and
Texture Entropy) were identified by substring matching on CellProfiler
feature names. Corrected direction instability was computed on the remaining
2{,}437 non-health features.

All scoring functions and curated drug lists were committed before any real
LINCS signatures were analyzed. Perturb-seq hypotheses (HP1--HP3) were
pre-registered before the correction analysis was run; the integrity
boundary, decision criteria, and commit SHA are documented in the
repository.

\paragraph{Test family and multiplicity.}
We treat the five LINCS drug hypotheses (H1--H5) as one test family and the
three Perturb-seq hypotheses (HP1--HP3) as a second. Within each family,
hypotheses test distinct validity conditions rather than the same quantity at
multiple thresholds, so we pre-registered individual decision criteria
($\alpha = 0.0125$ for HP1--HP3; specific thresholds for H1--H5) rather than
a family-wise correction. The pre-registered success criterion---3 of 5 drug
hypotheses confirmed---accounts for multiplicity at the study level.
After applying Holm--Bonferroni correction within the drug family, H3
($p < 10^{-27}$) and H5 (5/5 folds, binomial $p = 0.031$) survive at
$\alpha = 0.05$; H2 ($p = 0.28$, binomial 7/10) and H4 (gap $= 0.018$)
do not. H1's criterion-pass is an artifact (see Results). Within Perturb-seq,
HP3 ($p < 10^{-30}$) survives; HP1 and HP2 are characterization tests
(ranking correlations) without formal hypothesis tests.

\subsection{Toxicity-corrected bracket (Experiment 1)}

We curated 16 known cytotoxic compounds (staurosporine, camptothecin,
doxorubicin, etoposide, cisplatin, vincristine, paclitaxel, actinomycin-d,
thapsigargin, tunicamycin, brefeldin-a, rotenone, antimycin-a, oligomycin-a,
menadione, hydrogen-peroxide) and 41 stress/apoptosis marker genes
(\emph{HSPA1A}, \emph{BAX}, \emph{TP53}, \emph{CASP3}, \emph{ATF4},
\emph{DDIT3}, etc.). Toxicity-corrected bracket removes these gene axes from
the signature matrix before recomputing direction instability on the residual.

\textbf{Pre-registered prediction (H1):} $\geq$60\% of cytotoxic drugs move
from the top 20\% of raw bracket to below-median toxicity-corrected bracket.

\subsection{Broad-mechanism therapeutic bracket (Experiment 2)}

We curated 10 broad-mechanism drugs: statins (atorvastatin, simvastatin,
lovastatin), metformin, dexamethasone, sirolimus, methotrexate, aspirin,
and pan-HDAC inhibitors (vorinostat, trichostatin-a). We also curated 8
context-specific drugs (vemurafenib, imatinib, erlotinib, crizotinib,
lapatinib, gefitinib, PCI-34051, tubastatin-a).

\textbf{Pre-registered prediction (H2):} $\geq$60\% of broad-mechanism drugs
fall below the population median direction instability. Context-specific
drugs fall above the median.

\subsection{Phenotype-projected bracket (Experiment 3)}

Rung~5 of the validity ladder requires that direction instability align with
a known target pathway. We test this using genetic perturbation connectivity:
for each drug with an annotated target gene (Drug Repurposing Hub,
1{,}551 drugs with target annotations), the ``on-target direction'' is the
consensus shRNA knockdown signature of that target gene in LINCS L1000.
The shRNA signatures span 20 cell lines, dominated by a core set
(PC3, VCAP, MCF7, A375, HA1E, A549, HT29, HCC515, HEPG2; 9 lines
contributing $>$95\% of signatures). Each consensus averages across all
cell lines and hairpins for genes with $\geq$3 shRNA constructs
(4{,}369 target genes eligible).

Phenotype-projected bracket projects the instability vector onto this
on-target direction. On-target enrichment is quantified by cosine squared
between each drug's mean signature and its target gene's shRNA consensus
direction---the fraction of mean-signature variance aligned with the
genetic perturbation.

\textbf{Pre-registered prediction (H3):} Spearman $|\rho| > 0.3$ between
phenotype-projected bracket and on-target enrichment, AND $|\rho| < 0.15$
for raw bracket versus the same enrichment.

\subsection{Localization score predicts mechanism (Experiment 4)}

Rung~4 requires that instability concentrate in pathway-relevant genes.
For each drug, the region mask comprises the 100 genes with largest absolute
$z$-score in the target gene's shRNA consensus signature ($\sim$10\% of the
978 landmark genes). The localization score is the ratio of pathway-restricted
bracket to global bracket. MOA prediction uses macro-averaged one-vs-rest
AUROC across MOA classes with $\geq$10 drugs (15 classes, 223 eligible drugs).

\textbf{Pre-registered prediction (H4):} Macro-averaged AUROC for MOA
prediction from localization score $>$ AUROC from raw bracket by at least
0.05.

\subsection{Transport-stable bracket (Experiment 5)}

Rung~6 requires that the score replicate across held-out contexts.
Transport-stable bracket penalizes raw direction instability by Fr\'{e}chet
variance across contexts. We evaluate via leave-one-cell-line-out
cross-validation over the first 5 cell lines alphabetically (A375, A549,
A673, AGS, ASC). For each fold, we compute raw and transport-stable bracket
on the training cell lines and measure cosine similarity between the
held-out cell line's signature and the training consensus direction. This
held-out outcome matches the direction-based estimand of direction
instability: a drug whose signatures are directionally consistent in
training should produce a held-out signature aligned with the training
consensus.

\textbf{Pre-registered prediction (H5):} Spearman correlation between
transport-stable bracket and held-out cosine $>$ correlation for raw bracket,
in at least 4 of 5 folds.

\subsection{Essential-gene correction on Perturb-seq (Experiment 7)}

Universal-essential knockdowns (DepMap Chronos $< -0.5$ in both K562 and
RPE1) may inflate or deflate Grassmannian geodesic distances through a
shared lethal-response subspace. We computed the Fr\'echet mean subspace of
the 855 essential-gene perturbation responses (separately in each cell type)
via iterative projection on Gr$(5, 2073)$. We then projected this
essential-response subspace out of each gene's response subspace,
re-orthogonalized, and recomputed geodesic distances on the residual.

\textbf{Pre-registered prediction (HP1):} The essential-response correction
changes the geodesic-distance ranking. Decision criterion: Spearman $\rho$
between raw and corrected distances below 0.70 = ``bites''; 0.70--0.83 =
``marginal''; $\geq 0.83$ = ``doesn't bite'' (the HDAC selectivity
gradient gives $\rho = 0.83$ after drug toxicity correction).

\textbf{Pre-registered characterization (HP3):} Universal-essential
knockdowns produce higher geodesic distances than non-essential knockdowns
(Cohen's $d > 0$), consistent with the failure-mode inversion identified
before freezing.

\subsection{Pathway function prediction (HP2)}

If the essential-gene correction removes a confound that masks biological
signal, corrected distances should better predict Gene Ontology pathway
coherence. For 50{,}000 sampled gene pairs (from 1{,}544 GO-annotated genes
in our dataset), we computed pairwise Jaccard overlap of GO Biological
Process memberships (7{,}647 pathways from MSigDB C5, release 2023.2) and
correlated this with both the absolute distance difference $|d_i - d_j|$
and the mean distance $(d_i + d_j)/2$ for raw and corrected distances.

\textbf{Pre-registered prediction (HP2):} Spearman $|\rho|$ between corrected
distance and GO pathway Jaccard $>$ $|\rho|$ for raw distance, by at least
0.05.

\subsection{Cell-health correction on JUMP-CP (Experiment 8)}

Cell Painting morphological profiles include features that reflect general
cell health (area, compactness, DNA and mitochondrial intensity, granularity)
alongside compound-specific morphological effects (texture, spatial
distribution patterns in ER, AGP, RNA channels). Compounds that reduce
viability may achieve low direction instability through consistent
health-feature signatures rather than mechanism-specific morphology. We
identified 743 cell-health features by substring matching and recomputed
direction instability on the remaining 2{,}437 features for each of 25{,}254
compounds tested in $\geq$5 imaging sources.

\textbf{Pre-registered prediction:} Spearman $\rho$ between raw and
cell-health-corrected direction instability below 0.70 = ``bites'';
0.70--0.83 = ``marginal''; $\geq 0.83$ = ``doesn't bite'' (same decision
criteria as Experiments 1 and 7).

\section{Results}

\subsection{Experiment 1: Toxicity correction has negligible effect in LINCS}

Among the 16 pre-registered cytotoxic compounds, 13 were present in the
$\geq$5-cell-line dataset. Their raw direction instabilities ranged from
0.47 to 0.93 (Table~\ref{tab:toxic}). This immediately contradicts the
predicted failure mode: cytotoxic drugs are not false-positive transporters.
They are already high-instability---their signatures rotate across cell lines
rather than conserving a consistent stress direction.

After removing stress/apoptosis gene axes, direction instability shifted by
$<0.01$ for all 13 drugs. The correction has no practical effect because the
premise of the failure mode---that cytotoxic drugs achieve low bracket via
a shared stress program---does not hold in LINCS L1000 data.

\begin{table}[t]
\centering
\caption{Direction instability for known cytotoxic drugs. Most are above the
population median (0.92), contradicting the predicted failure mode.
Correction for stress genes shifts values by $<$0.01.}
\label{tab:toxic}
\small
\begin{tabular}{lrrr}
\toprule
Drug & $D_{\text{raw}}$ & $D_{\text{corr}}$ & $\Delta$ \\
\midrule
staurosporine & 0.47 & 0.48 & $+$0.01 \\
doxorubicin & 0.74 & 0.74 & 0.00 \\
etoposide & 0.78 & 0.78 & 0.00 \\
camptothecin & 0.82 & 0.82 & 0.00 \\
vincristine & 0.87 & 0.87 & 0.00 \\
paclitaxel & 0.89 & 0.89 & 0.00 \\
cisplatin & 0.91 & 0.91 & 0.00 \\
tunicamycin & 0.92 & 0.92 & 0.00 \\
thapsigargin & 0.93 & 0.93 & 0.00 \\
\midrule
\emph{Population median} & 0.92 & --- & --- \\
\bottomrule
\end{tabular}
\end{table}

\paragraph{Interpretation.}
H1 passes its pre-registered criterion (85\% of toxic drugs fall below
the corrected median), but this is an artifact of the criterion design: most
toxic drugs are already high-instability, so correction moves them
negligibly rather than reclassifying them. We classify H1 as
\emph{refuted with artifact pass}: the predicted mechanism (cytotoxic
drugs masquerading as mechanism-conserving via shared stress programs) is
wrong. Cytotoxicity produces genuinely context-dependent signatures.
The failure mode---``violence mimics transport''---is not the dominant
failure mode in 978-gene pharmacological data.

\subsection{Experiment 2: Broad-mechanism drugs confirm; specificity gradient is weaker than predicted}

Seven of 10 broad-mechanism drugs fell below the population median direction
instability of 0.92, confirming H2 (Table~\ref{tab:broad}). The strongest
signal came from pan-HDAC inhibitors: vorinostat ($D = 0.49$) and
trichostatin-a ($D = 0.55$) are among the most mechanism-conserving drugs
in the entire dataset.

\begin{table}[t]
\centering
\caption{Direction instability for pre-registered broad-mechanism drugs.
7/10 fall below the population median (0.92). Pan-HDAC inhibitors dominate
the low end. Statins and aspirin are above median---``universal target'' does
not guarantee a conserved transcriptomic signature.}
\label{tab:broad}
\small
\begin{tabular}{llrr}
\toprule
Drug & Rationale & $D$ & Below median? \\
\midrule
vorinostat & pan-HDAC & 0.49 & Yes \\
trichostatin-a & pan-HDAC & 0.55 & Yes \\
sirolimus & mTOR & 0.78 & Yes \\
methotrexate & DHFR & 0.81 & Yes \\
dexamethasone & GR agonist & 0.83 & Yes \\
metformin & AMPK & 0.86 & Yes \\
lovastatin & HMGCR & 0.88 & Yes \\
atorvastatin & HMGCR & 0.91 & No \\
simvastatin & HMGCR & 0.93 & No \\
aspirin & COX & 0.96 & No \\
\bottomrule
\end{tabular}
\end{table}

\paragraph{The unexpected kinase result.}
The pre-registered second prediction of H2---that context-specific kinase
inhibitors would fall above the median---failed. All 8 tested kinase
inhibitors (vemurafenib, imatinib, erlotinib, crizotinib, lapatinib,
gefitinib, PCI-34051, tubastatin-a) had direction instability below the
population median, with values ranging from 0.73 to 0.91
(Table~\ref{tab:specific}).

\begin{table}[t]
\centering
\caption{Direction instability for pre-registered context-specific drugs.
Contrary to prediction, all fall below the population median.}
\label{tab:specific}
\small
\begin{tabular}{llrr}
\toprule
Drug & Target & $D$ & Above median? \\
\midrule
vemurafenib & BRAF V600E & 0.73 & No \\
imatinib & BCR-ABL & 0.78 & No \\
erlotinib & EGFR & 0.81 & No \\
crizotinib & ALK/ROS1 & 0.83 & No \\
lapatinib & HER2 & 0.84 & No \\
gefitinib & EGFR & 0.86 & No \\
PCI-34051 & HDAC8 & 0.89 & No \\
tubastatin-a & HDAC6 & 0.91 & No \\
\bottomrule
\end{tabular}
\end{table}

\paragraph{Interpretation.}
The sorting axis is target \emph{breadth}---how many gene programs a compound
engages---not target \emph{specificity}. Pan-HDAC inhibitors rank lowest
because chromatin remodeling affects hundreds of downstream genes
simultaneously, producing a strong, coherent direction. Kinase inhibitors
target fewer downstream genes but still target them consistently across cell
types (EGFR signaling is active in most LINCS cell lines, regardless of
whether EGFR mutation drives growth). The drugs that actually have high
direction instability ($D > 0.94$) tend to be receptor agonists and
neuromodulators whose targets are tissue-restricted in expression.

The gradient we observe is: conserved engagement of many genes (low $D$)
$\to$ conserved engagement of few genes (moderate $D$) $\to$ tissue-restricted
expression of target (high $D$). Context-specificity at the clinical level
(vemurafenib only works in BRAF-mutant melanoma) does not imply context-specificity
at the transcriptomic level (vemurafenib still engages MAPK signaling in
non-mutant lines, just without therapeutic consequence).

\subsection{Experiment 3: Phenotype-projected bracket separates on-target from off-target instability}

Among 795 drugs with annotated targets and matched shRNA consensus signatures,
phenotype-projected bracket correlates with on-target enrichment at Spearman
$\rho = 0.376$ ($p = 4.8 \times 10^{-28}$). Raw bracket shows no correlation
with the same enrichment ($\rho = -0.043$, $p = 0.22$). Both pre-registered
criteria are satisfied: $|\rho_{\text{projected}}| = 0.376 > 0.3$ and
$|\rho_{\text{raw}}| = 0.043 < 0.15$.

\paragraph{Interpretation.}
H3 confirms Rung~5 of the validity ladder empirically. Raw direction
instability measures whether a drug's signatures point in consistent
directions across cell lines, but this consistency could arise from any
coherent program---on-target, off-target, or stress-related. Projecting
instability onto the genetic perturbation direction of the drug's annotated
target isolates the mechanistically relevant component. The near-zero
correlation for raw bracket shows that global instability carries no
information about on-target alignment; the strong correlation for projected
bracket shows that the on-target projection recovers it.

The gene-set construction avoids researcher degrees of freedom: the shRNA
consensus signature is an empirical object from the same LINCS L1000
platform, requiring no pathway database selection or gene-set size
decisions. H3 and H4 share this shRNA dependency
(documented in DEVIATION\_LOG.md entry~3) and are not independent tests.

\subsection{Experiment 4: Localization predicts intracellular but not receptor-mediated mechanism}

\paragraph{H4 is not confirmed.}
The macro-averaged one-vs-rest AUROC gap between localization score and raw
bracket is $+0.018$ (95\% bootstrap CI $[-0.010, +0.048]$), below the
pre-registered threshold of $+0.05$. The confidence interval includes zero.

\paragraph{The failure has biological structure.}
Per-MOA AUROC gaps reveal a consistent pattern: localization improves
prediction for drugs targeting intracellular enzymes and diverges for
receptor-mediated drugs (Table~\ref{tab:loc_moa}).

\begin{table}[t]
\centering
\caption{Per-MOA AUROC gaps (localization $-$ raw) for the 15 eligible MOA
classes. Intracellular enzyme targets (topoisomerase, CDK, mTOR, p38~MAPK)
show large positive gaps. Receptor-mediated targets (acetylcholine, PPAR,
estrogen, histamine) show negative gaps. Macro-average: $+0.018$.}
\label{tab:loc_moa}
\small
\begin{tabular}{lrr}
\toprule
MOA class & $N$ & AUROC gap \\
\midrule
topoisomerase inhibitor & 15 & $+$0.491 \\
CDK inhibitor & 16 & $+$0.258 \\
mTOR inhibitor & 13 & $+$0.176 \\
p38 MAPK inhibitor & 10 & $+$0.165 \\
HDAC inhibitor & 20 & $+$0.076 \\
glucocorticoid receptor agonist & 27 & $+$0.061 \\
MEK inhibitor & 13 & $+$0.036 \\
\midrule
adrenergic receptor agonist & 14 & $-$0.036 \\
calcium channel blocker & 11 & $-$0.087 \\
cyclooxygenase inhibitor & 17 & $-$0.099 \\
EGFR inhibitor & 39 & $-$0.107 \\
histamine receptor antagonist & 26 & $-$0.135 \\
estrogen receptor agonist & 16 & $-$0.148 \\
PPAR receptor agonist & 16 & $-$0.188 \\
acetylcholine receptor antagonist & 23 & $-$0.201 \\
\bottomrule
\end{tabular}
\end{table}

\paragraph{Interpretation.}
Localization to the shRNA-derived gene set works when the drug's mechanism
operates through a small number of genes that the target gene's knockdown
also affects---topoisomerase inhibitors, CDK inhibitors, mTOR inhibitors
produce transcriptomic effects concentrated in the same gene programs as
their target knockdowns. Receptor-mediated drugs produce diffuse downstream
cascades that spread beyond the 100-gene region mask. Their instability is
real but not localizable to the target gene's knockdown signature.

This is an informative null. The localization rung of the validity ladder
is domain-dependent: it adds inferential warrant for intracellular enzyme
inhibitors but not for receptor agonists and antagonists, whose mechanisms
produce distributed transcriptomic effects. The pre-registered threshold
appropriately penalizes this heterogeneity through the macro-average.

\subsection{Experiment 5: Transport-stable bracket outpredicts raw bracket in held-out cell lines}

Transport-stable bracket outpredicts raw bracket in all 5 cross-validation
folds (Table~\ref{tab:transport}). Raw bracket is \emph{anti-predictive}:
its Spearman correlation with held-out cosine consistency is negative in
every fold, reaching $\rho = -0.73$ in the A673 fold. Transport-stable
bracket is positively predictive in all folds, reaching $\rho = +0.72$ in
the same fold. The pre-registered criterion (transport-stable wins in
$\geq$4 of 5 folds) is satisfied at 5/5.

\begin{table}[t]
\centering
\caption{Leave-one-cell-line-out cross-validation. Transport-stable (TS)
bracket outpredicts raw bracket in all 5 folds. Raw bracket is
anti-predictive: higher raw instability in training predicts \emph{lower}
held-out cosine consistency, while higher TS bracket predicts higher
consistency.}
\label{tab:transport}
\small
\begin{tabular}{lrrrl}
\toprule
Held-out cell & $N_{\text{drugs}}$ & $\rho_{\text{raw}}$ & $\rho_{\text{TS}}$ & Winner \\
\midrule
A375 & 8{,}328 & $-$0.480 & $-$0.155 & TS \\
A549 & 8{,}888 & $-$0.494 & $-$0.201 & TS \\
A673 & 357 & $-$0.725 & $+$0.725 & TS \\
AGS & 356 & $-$0.712 & $+$0.714 & TS \\
ASC & 2{,}345 & $-$0.365 & $+$0.346 & TS \\
\bottomrule
\end{tabular}
\end{table}

\paragraph{The effect is strongest in rare cell lines.}
Folds holding out rare cell lines (A673: 357 drugs, AGS: 356 drugs) show
the largest absolute correlations and the starkest sign flips. In these
folds, removing a single cell line from training substantially changes the
consensus direction, and drugs with high raw instability (inconsistent
directions in training) produce held-out signatures that point away from
the training consensus. Transport-stable bracket corrects for this by
penalizing high Fr\'{e}chet variance, so drugs whose training instability
is driven by a single outlier cell line are downweighted.

\paragraph{Common cell lines show the same direction, smaller magnitude.}
In the A375 and A549 folds (8{,}000+ drugs), raw bracket remains
anti-predictive ($\rho \approx -0.49$) and transport-stable bracket
improves over raw, though both remain negative. The consensus direction
is stable when computed from many cell lines, so the held-out prediction
task is easier and both metrics perform more similarly---but the ordering
is preserved.

\paragraph{Interpretation.}
H5 provides the strongest single result in the drug domain. Raw direction
instability is anti-predictive for cross-context transport: drugs with high
instability in training produce held-out signatures that disagree with the
training consensus. This occurs because high instability conflates two
distinct phenomena---genuine directional inconsistency and Fr\'{e}chet
variance (sampling noise from a small number of contexts). Transport-stable
bracket separates these, and the separation matters most precisely where the
confound is strongest: in rare cell lines where a single outlier can
dominate the raw metric.

\paragraph{Shared-variance caveat.}
The transport-stable penalty and the held-out cosine outcome both depend on
the training-set signature geometry, raising a potential circularity concern:
both quantities share variance through the training consensus direction.
The sign flip (raw anti-predictive, TS predictive) argues against a trivial
shared-variance explanation---if the improvement were purely algebraic,
we would expect monotonic improvement rather than a sign reversal. A
permutation baseline (shuffling drug labels within each fold) would
provide a formal null; we leave this for future work.

\subsection{Experiment 7: Essential-gene correction preserves ranking but reveals confound-dominated machinery}

\paragraph{HP3: Failure-mode inversion.}
Universal-essential knockdowns have \emph{higher} Grassmannian geodesic
distances than non-essential knockdowns (mean 2.80 vs.\ 2.68, Cohen's $d =
+0.55$, 95\% CI $[0.45, 0.65]$, Mann--Whitney $p < 10^{-30}$). Essential
genes do not produce a shared response across cell types; they
\emph{diverge}---the same gene's lethal disruption manifests through
different transcriptomic programs in K562 and RPE1. This inverts the
predicted failure mode: in drugs, ``violence'' (cytotoxicity) was predicted
to mimic transport but does not; in genetics, ``violence'' (essential
knockdowns) produces the opposite of transport. Because both screens target
DepMap-essential gene libraries (85\% overlap), the ``non-essential'' group
still contains moderately essential genes, making the observed $d = 0.55$
effect a conservative estimate of the true essential/non-essential divergence.

\paragraph{HP1: Correction does not bite.}
After projecting out the essential-response Fr\'echet mean subspace,
Spearman $\rho$ between raw and corrected geodesic distances is 0.91 ($p <
10^{-300}$). By the pre-registered criterion, this ``doesn't bite''---the
same qualitative outcome as the drug HDAC gradient ($\rho = 0.83$).
The global ranking is preserved.

\paragraph{Item-level rank shifts are the primary finding.}
The essential-gene correction's value is not at the population level
(where $\rho = 0.91$ confirms stability) but at the item level: individual
genes shift hundreds of rank positions, identifying which transport scores
are most confound-dominated. The top 100 rank-shifters are enriched for
spliceosomal complex assembly (Fisher's exact OR $= 8.9$, $p = 1.3 \times
10^{-7}$) and proteasomal catabolic processes (OR $= 4.1$, $p = 6.7 \times
10^{-5}$): proteasome subunits (\emph{PSMA3}, \emph{PSMB6}, \emph{PSMC4})
shift 750--900 positions; spliceosome components (\emph{SF3B1},
\emph{SNRPG}, \emph{DDX23}) shift 660--820 positions. These core-machinery
complexes are the genes whose cross-cell-type geometry is most dominated by
the shared essential-response axis. The correction provides a concrete
action: flag these genes before interpreting their transport scores as
mechanistically meaningful.

\paragraph{HP2: Pathway function prediction does not bite.}
GO Biological Process Jaccard overlap has near-zero correlation with geodesic
distance in both raw ($\rho = 0.024$, $p < 10^{-7}$) and corrected ($\rho =
-0.002$, $p = 0.68$) forms. The correction removes the tiny positive
relationship that existed rather than strengthening it (improvement $=
-0.022$). By the pre-registered criterion (improvement $\geq 0.05$), HP2
does not bite. Geodesic distance measures transcriptomic-geometry rotation,
which is orthogonal to pathway co-membership at this scale: two genes in the
same GO category need not produce geometrically similar perturbation
responses.

\paragraph{Interpretation.}
The Perturb-seq domain provides the same qualitative result as LINCS drugs:
the validity condition (essential-gene correction) does not reorganize the
global ranking but upgrades the interpretive warrant. The domain-specific
insight is the failure-mode inversion: universal-essential knockdowns
diverge rather than converge, which means the confounding mechanism operates
through a different channel than in pharmacological perturbations. Cohen's $d$
drops from $+0.55$ to $+0.28$ after correction, confirming that the
essential-response subspace accounts for roughly half the group separation.
HP2's null further constrains interpretation: geodesic distance captures
geometric similarity of perturbation effects, which is distinct from
functional pathway relatedness.

\subsection{Experiment 8: Cell-health correction has negligible effect in JUMP-CP}

Among 25{,}254 compounds tested in $\geq$5 imaging sources, raw direction
instability has population mean 0.92 and median 0.95---comparable to the
LINCS drug population (median 0.92 in 978-gene space). After removing 743
cell-health features, corrected direction instability tracks the raw ranking
almost perfectly: Spearman $\rho = 0.987$ ($p < 10^{-300}$). By the
pre-registered criterion, this ``doesn't bite.''

\paragraph{The correction is even weaker than in the other domains.}
The JUMP-CP $\rho = 0.987$ is comparable to the LINCS global drug
correction ($\rho > 0.99$) and slightly above the Perturb-seq
essential-gene correction ($\rho = 0.91$). This
is consistent with the higher feature dimensionality: 743 health features
constitute 23\% of the 3{,}180-feature space, a smaller proportional
perturbation than the drug toxicity gene set (41 of 978 genes, 4\%) yet
with less ranking impact, because the health features carry less of the
cross-source variance in morphological signatures.

\paragraph{Individual compounds still shift thousands of ranks.}
Despite the near-perfect global correlation, individual compounds shift up
to 7{,}155 rank positions after cell-health correction
(Table~\ref{tab:jump_movers}). These are compounds whose direction
instability is dominated by cell-health feature variance---their
morphological signatures appear consistent across sources primarily because
health-related measurements (area, granularity, mitochondrial intensity)
behave similarly, rather than because the compound produces a
mechanism-specific morphological phenotype.

\begin{table}[t]
\centering
\caption{Top rank movers after cell-health correction in JUMP-CP. Despite
$\rho = 0.987$ globally, individual compounds shift thousands of positions,
identifying those whose apparent morphological consistency is
health-feature-dominated.}
\label{tab:jump_movers}
\small
\begin{tabular}{lrrr}
\toprule
Compound & Raw rank & Corrected rank & $|\Delta|$ \\
\midrule
JCP2022\_015169 & 10{,}843 & 17{,}998 & 7{,}155 \\
JCP2022\_098151 & 16{,}600 & 23{,}618 & 7{,}018 \\
JCP2022\_093366 & 8{,}306 & 15{,}020 & 6{,}714 \\
JCP2022\_102324 & 14{,}753 & 21{,}442 & 6{,}689 \\
JCP2022\_037124 & 9{,}200 & 15{,}800 & 6{,}600 \\
\bottomrule
\end{tabular}
\end{table}

\paragraph{Interpretation.}
JUMP-CP provides the strongest test of the ``correction doesn't bite''
pattern. With 25{,}254 compounds, 3{,}180 features, and $K \geq 5$ imaging
sources providing genuine direction-instability computation (not the $K = 2$
pairwise reduction of Perturb-seq), the cell-health correction yields $\rho
= 0.987$. The morphological domain joins the transcriptomic (LINCS) and
genetic (Perturb-seq) domains in exhibiting the same pattern: validity
conditions preserve global rankings while identifying individual
perturbations whose scores are most confound-dominated.

\section{Discussion}

\subsection{Validity conditions license claims, they do not reclassify}

The central finding is negative in the most useful sense: validity conditions
do not reorganize the direction-instability ranking. In LINCS drugs, the
HDAC selectivity gradient---the fine-grained ordering from pan-isoform to
selective inhibitors---is preserved after toxicity correction ($\rho =
0.83$; the global 8{,}949-drug ranking gives $\rho > 0.99$). In Perturb-seq
knockdowns, the essential-gene correction preserves the global
geodesic-distance ranking ($\rho = 0.91$). In JUMP-CP, the cell-health
correction leaves the morphological ranking nearly unchanged ($\rho =
0.99$). All three domains tell the same story: the ranking is robust.

What changes is the inferential warrant. Phenotype projection (H3) proves
that direction instability can be decomposed into on-target and off-target
components: the projected bracket correlates with genetic perturbation
connectivity ($\rho = 0.38$) while the raw bracket carries no on-target
information. Transport stability (H5) proves that the metric predicts
held-out cell-line effects, but only after Fr\'{e}chet variance correction;
raw bracket is anti-predictive, conflating genuine directional inconsistency
with sampling noise. Applying toxicity correction and
observing negligible movement proves that the HDAC signal is
chromatin-specific, not stress-derived. Applying essential-gene correction
and observing that proteasome/spliceosome genes shift while the population
ranking holds proves that the transport signal is not an artifact of shared
lethal-response geometry.

Confound corrections for perturbation scores are well-established
\citep{meyers2017ceres,pacini2021integrated}; the novelty here is not that
corrections exist but that rank-preserving corrections carry inferential
value. The validity ladder is a proof system: each rung eliminates a specific
threat to interpretation. A practitioner who reports only raw $D = 0.55$
makes a weaker claim than one who reports $D = 0.55$ after toxicity
correction, localization to the HDAC pathway, and replication across a
held-out cell-line panel. The number is the same; the license differs.
This aligns with recent work on construct validity for computational
scores \citep{freiesleben2025benchmarking,watson2025measuring}: the
question is not whether the metric is accurate but what inferences it
supports.

\subsection{Honest nulls strengthen the contribution}

Five of our pre-registered predictions produced informative nulls or refutations.

\paragraph{H1 is refuted; its criterion-pass is an artifact.}
We predicted that cytotoxic drugs would masquerade as mechanism-conserving
(low bracket) via shared stress programs, and that toxicity correction would
expose them. In LINCS L1000, cytotoxic drugs are already high-instability.
Their transcriptomic effects are genuinely context-dependent---different stress
cascades dominate in different cell types. H1's pre-registered criterion
(85\% of toxic drugs below corrected median) passes because most toxic drugs
are high-instability to begin with, not because the correction reclassifies
them. We classify this as \emph{refuted with artifact pass}: the predicted
mechanism is wrong, and the criterion is too permissive to detect this.
The ``violence mimics transport'' failure mode may be more relevant in
lower-dimensional assays where the stress response dominates the
measurable signal.

\paragraph{H2's specificity axis was wrong.}
We predicted that context-specific kinase inhibitors would have high direction
instability. They do not. The sorting axis is target breadth (number of
affected gene programs), not clinical specificity. This distinguishes the
transcriptomic notion of ``context-dependent effect'' from the clinical notion
of ``specific to patient subgroup.'' Both are informative; they are not the
same quantity.

\paragraph{H4's localization is target-class-dependent.}
We predicted that localization to target-gene knockdown signatures would
improve MOA prediction by at least 0.05 AUROC. The macro-averaged gap is
$+0.018$. Localization helps substantially for intracellular enzyme
inhibitors (topoisomerase: $+0.49$, CDK: $+0.26$) but hurts for
receptor-mediated drugs (acetylcholine: $-0.20$, PPAR: $-0.19$). The
pre-registered macro-average threshold appropriately captures this
heterogeneity: Rung~4 is not universally applicable. For drugs whose
downstream effects are distributed across many gene programs, concentrating
the instability measurement on 100 target-gene-derived genes discards
informative signal.

\paragraph{HP1 doesn't bite.}
We predicted that essential-gene correction would disrupt the Perturb-seq
ranking. It does not ($\rho = 0.91$). This parallels the drug domain ($\rho
> 0.99$ globally): the correction preserves population-level rankings. The
failure-mode inversion
(HP3, essential knockdowns diverge rather than converge) explains why:
universal-essential genes do not create a shared response axis that inflates
transport scores, because their responses rotate \emph{away} from each other
across cell types.

\paragraph{HP2 doesn't bite.}
We predicted that corrected distances would better predict GO pathway
function. They do not---the relationship between geodesic distance and
pathway Jaccard overlap is negligible ($|\rho| < 0.025$) for both raw and
corrected distances. This separates two quantities that could have been
conflated: geometric similarity of perturbation responses is not the same as
functional pathway relatedness.

These nulls constrain the theory. Validity conditions are needed when the raw
metric's failure mode is active in a domain. In both LINCS drugs and
Perturb-seq genetics, the predicted confounding mechanism operates differently
than expected, and the correction adds little to the global ranking. The
ladder's value is domain-dependent, which is precisely the point: validity
conditions are not universal corrections but domain-specific proof obligations.

\subsection{Scope of claims}

Direction instability measures angular consistency of perturbation signatures
across contexts. It does not measure therapeutic efficacy, in-vivo
translatability, or clinical specificity. A drug with low direction
instability and full ladder validation has a strong claim to
\emph{mechanism-conserving transcriptomic effects across cell lines in
vitro}. Whether that mechanism is therapeutically relevant, whether the
effect replicates in patient-derived tissue, and whether the consistency
reflects target engagement rather than pathway compensation are questions
the ladder does not answer. The ladder is a proof system for
interpretive warrant, not a surrogate endpoint.

\subsection{Cross-domain consistency}

The same geometric object---angular deviation of perturbation-response
directions across contexts---captures mechanism conservation in three
unrelated domains: small-molecule drug effects across cell lines (LINCS
L1000, 978 genes, $K \geq 5$), genetic knockdown effects across
cell types (Perturb-seq, 2{,}073 genes, $K = 2$), and morphological
compound effects across imaging sources (JUMP-CP, 3{,}180 features,
$K \geq 5$). All three domains exhibit the same pattern: the relevant
validity condition preserves the global ranking while identifying specific
perturbations whose scores are most confound-dominated
(Table~\ref{tab:cross_domain}).

\begin{table}[t]
\centering
\caption{Cross-domain summary. The correction-vs-raw Spearman $\rho$
quantifies how much the domain-specific validity condition reorganizes
the ranking. All three domains exceed the pre-registered ``doesn't bite''
threshold ($\rho \geq 0.83$). The LINCS HDAC gradient $\rho$ measures
preservation of the fine-grained selectivity ordering within a single drug
class (20 compounds); the global 8{,}949-drug $\rho > 0.99$.}
\label{tab:cross_domain}
\small
\begin{tabular}{llrrl}
\toprule
Domain & Correction & Features & $N$ & $\rho$ \\
\midrule
LINCS L1000 (HDAC gradient) & Toxicity genes & 978 & 20 & 0.83 \\
LINCS L1000 (global) & Toxicity genes & 978 & 8{,}949 & $>$0.99 \\
Perturb-seq & Essential-gene subspace & 2{,}073 & 1{,}676 & 0.91 \\
JUMP-CP & Cell-health features & 3{,}180 & 25{,}254 & 0.99 \\
\bottomrule
\end{tabular}
\end{table}

The domain-specific failure modes differ in mechanism: in drugs, cytotoxicity
produces high-instability signatures (context-dependent stress cascades); in
genetics, universal-essential knockdowns produce high-distance subspaces
(divergent lethal programs); in morphology, cell-health features contribute
a small fraction of the cross-source variance. All three invert the naive
prediction that ``violence mimics transport.'' The validity ladder provides
the domain-specific proof structure needed to interpret a general measurement
within each domain's causal framework.

\paragraph{Robustness extends to process-view choice.}
To test whether restricting contexts to a single tissue type would create
enough confounding for correction to bite, we computed within-tissue
direction instability for drugs tested in $\geq$5 cell lines of a single
tissue. Across five tissues (large intestine, lung, heme, breast, ovary;
176--358 drugs each), lineage-gene correction (removing the top-100
within-tissue variance genes) yields $\rho = 0.98$--$0.99$ in all cases.
The ranking is robust even when contexts share tissue identity. This
result extends the pattern: in high-dimensional feature spaces, removing
a biologically motivated feature subset (whether toxicity genes, essential
response, cell-health features, or lineage genes) does not reorganize the
global ranking.

The one context where correction has measurable effect is within-class
comparisons at low $N$: the 20-compound HDAC selectivity gradient gives
$\rho = 0.83$, and within-essential-gene ranking in Perturb-seq gives $\rho
= 0.88$. This suggests the ladder's empirical bite is concentrated in
fine-grained mechanistic comparisons within a single perturbation class,
where the confound can dominate the limited between-perturbation variance.

\subsection{Relation to the drug transport paper}

\citet{tower2026bracket} established that direction instability predicts mechanism
conservation and demonstrated three internal validity checks: the cytotoxicity
defense (ruling out stress as a driver), the transporting core (identifying
conserved genes), and genetic triangulation (confirming target-mediated
transport via shRNA knockdown). Those analyses address validity implicitly---each
rules out a specific confound without naming the general principle.

This paper makes the general principle explicit and tests each rung
empirically. The cytotoxicity defense is
Rung~4 (localized to biology) applied negatively: showing that the signal is
\emph{not} in stress genes. The transporting core is Rung~5 (phenotype-aligned):
the conserved direction matches the target pathway---confirmed here via
genetic perturbation connectivity (H3, $\rho = 0.38$). Genetic triangulation is
Rung~7 (predicts holdout): an independent data modality confirms the mechanistic
prediction---extended here via cross-context transport stability (H5, 5/5
folds). Naming the ladder lets practitioners apply these proof obligations
to new domains (single-cell perturbation, genomic screens, morphological
assays) where the specific checks from the drug transport paper do not
directly translate.

\subsection{Limitations}

We tested five drug-domain experiments (H1--H5), three genetic-domain
experiments (HP1--HP3), and one morphological experiment (cell-health
correction). The drug-domain experiments H3--H5 use shRNA knockdown
signatures as the ground-truth on-target direction, which assumes the
target gene annotation in the Drug Repurposing Hub is correct and that
the shRNA consensus captures the drug's mechanism. Drugs with
unannotated or incorrectly annotated targets are excluded (7{,}398 of
8{,}949 drugs lack target annotations), and multi-target drugs use only
the first listed target.

LINCS L1000 measures 978 landmark genes, restricting the toxicity gene set to
those represented in the landmark panel. A full-transcriptome dataset would
test whether the negligible effect of toxicity correction persists at higher
dimensionality, where stress programs span more measurable axes.

The Perturb-seq analysis uses $K = 2$ cell types, which collapses the
Grassmannian analysis to a single pairwise comparison. There is no Fr\'echet
variance, no holonomy, no cross-context transport stability test. The
essential-gene correction result should be interpreted as a
correction-changes-ranking test rather than a full transport stability
analysis. LINCS and JUMP-CP both use $K \geq 5$ contexts, providing
genuine direction-instability computation in two of the three domains.
Both the K562 and RPE1 screens target DepMap-essential genes, so the
HP3 ``non-essential'' comparison group consists of genes selected for
essentiality that fall below the Chronos $< -0.5$ threshold in one or
both cell types. This attenuates the essential/non-essential contrast and
makes the observed effect conservative; a genome-wide Perturb-seq screen
with genuine non-essential controls would provide a sharper test.

JUMP-CP profiles are variance-stabilized and MAD-normalized, and the
cell-health feature set was defined by substring matching on CellProfiler
feature names rather than by an independent biological criterion. A more
principled partition (e.g., features correlated with viability assays)
could yield different correction magnitudes.

\subsection{Implications}

The validity ladder has immediate practical utility in three settings
demonstrated here:

\textbf{Drug repurposing.} A compound with low direction instability and
confirmed localization to a disease-relevant pathway has stronger repurposing
evidence than one with low instability alone. The ladder provides a checklist
for what additional evidence strengthens the repurposing case.

\textbf{Genetic screen prioritization.} Perturb-seq experiments produce
thousands of knockdown profiles. Grassmannian geodesic distance identifies
genes whose perturbation effects transport across cell types. The
essential-gene correction identifies which of these transport scores may be
confound-dominated, enabling more targeted follow-up: genes with low
geodesic distance \emph{after} correction are the strongest candidates for
conserved mechanism, while genes that shift substantially (proteasome,
spliceosome) warrant additional scrutiny before mechanistic claims.

\textbf{Morphological compound screening.} JUMP-CP direction instability
across 25{,}254 compounds provides a morphological analogue to the
transcriptomic ranking. The cell-health correction identifies compounds
whose apparent morphological consistency is dominated by viability-related
features rather than mechanism-specific phenotypes. For high-throughput
screening campaigns, this distinction matters: a compound that scores low
on direction instability \emph{after} cell-health correction has a stronger
claim to mechanism-conserving morphological effects than one whose
consistency derives primarily from cell-health signals.

\section{Conclusion}

Direction instability is a reliable measure of perturbation-response
consistency across contexts, whether applied to drug signatures in 978-gene
transcriptomic space, genetic knockdown responses as subspaces on the
Grassmannian, or compound morphological profiles in 3{,}180-feature Cell
Painting space. Validity conditions do not improve the metric's ability to
rank perturbations---the ranking is already biologically meaningful in all
three domains tested ($\rho > 0.99$ for LINCS, $0.91$ for Perturb-seq,
and $0.99$ for JUMP-CP after correction).
They improve the practitioner's ability to make claims about that ranking.
Phenotype projection separates on-target from off-target instability
($\rho = 0.38$ vs.\ $-0.04$); transport-stable bracket predicts held-out
effects where raw bracket is anti-predictive (5/5 folds); localization
helps for intracellular enzyme targets but not receptor-mediated drugs,
constraining where Rung~4 applies. Each condition eliminates a specific
interpretive threat: toxicity correction rules out stress-mediated false
positives; essential-gene correction rules out shared lethal-response
geometry; cell-health correction rules out viability-driven consistency.
The domain-specific failure modes differ, but the pattern is the same:
correction preserves population rankings while identifying individual
perturbations most susceptible to confounding. The validity ladder is a
proof system for mechanistic inference from geometric perturbation scores.

\bibliographystyle{tmlr}
\bibliography{references}

\end{document}
